% !TeX spellcheck = pl_PL-Polish
\documentclass[12pt,a4paper]{article}

\usepackage[T1]{fontenc}
\usepackage[polish]{babel}
\usepackage[utf8]{inputenc}
\usepackage{lmodern}
\selectlanguage{polish}
\usepackage{graphicx}
\usepackage{xcolor}
\usepackage{geometry}
\geometry{a4paper, margin=2cm}
\usepackage{hyperref}
\usepackage{listings}

\def\projectName{Zadanie 4 Quiz}
\def\authorA{Dawid Piotrowski (Generator quizu)}
\def\authorB{Filip Kłos (Rozwiązywanie quizu)}
\def\authorC{Mateusz Szymczak (Drugie rozwiązywanie quizu QuizSolver-M)}
\newcommand\todo[1]{\textcolor{red}{#1}}

% Ustawienia do wyświetlania kodu źródłowego
\lstset{
    language=[Sharp]C,
    backgroundcolor=\color{black!5},
    basicstyle=\ttfamily\footnotesize,
    keywordstyle=\color{blue},
    stringstyle=\color{red},
    commentstyle=\color{green!50!black},
    breaklines=true,
    captionpos=b,
    numbers=left,
    numberstyle=\tiny\color{gray},
    showstringspaces=false
}

\begin{document}
\pagenumbering{gobble}
\clearpage
	\begin{figure}[h]
		\centering
		%\includegraphics[width=0.5\linewidth]{media/polsl_logo_pion_en_rgb.png}
		\includegraphics[width=0.5\linewidth]{media/ps-logo.png}
	\end{figure}

\hspace{3cm}
	\begin{center}Dokumentacja projektowa\end{center}
	\hspace{3cm}
	\begin{center}\large\textbf{Zarządzanie Systemami Informatycznymi}\end{center}
	\begin{center}\large\textit{\projectName}\end{center}

\hspace{7cm}
	\begin{flushright}Kierunek: Informatyka
		\end{flushright}
		\begin{flushright}Członkowie zespołu:
		\par
		\textit{\authorA}
		\par
		\textit{\authorB}
		\par
		\textit{\authorC}
	\end{flushright}
\vfill
	\begin{center}Gliwice, 2025/2026\end{center}

\newpage
\pagenumbering{arabic}
\tableofcontents
\newpage

\section{Wprowadzenie}
\subsection{Role w projekcie}
Projekt został zrealizowany w zespole, gdzie podzielono się następującymi obowiązkami:
\begin{itemize}
    \item \textbf{\authorA:} Opracowanie aplikacji rozwiązującej quiz z wykorzystaniem architektury trójwarstwowej (wzorzec MVP) oraz interfejsu użytkownika w Windows Forms.
    \item \textbf{\authorB:} Opracowanie generatora quizu.
    \item \textbf{\authorC:} Opracowanie alternatywnego modułu rozwiązywania/rozszerzenia funkcjonalności.
\end{itemize}

\subsection{Cel projektu}
Celem projektu jest stworzenie aplikacji okienkowej składającej się z modułów: generatora quizu oraz aplikacji do rozwiązywania quizów. Aplikacja uczy budowy systemów w oparciu o architekturę trójwarstwową, w tym przypadku z wykorzystaniem wzorca projektowego Model-View-Presenter (MVP).

\section{Założenia projektowe}
\subsection{Założenia techniczne i nietechniczne}
\begin{itemize}
    \item Architektura trójwarstwowa (np. MVP).
    \item Pliki z quizami są szyfrowane (np. za pomocą algorytmu AES lub TripleDES) i deszyfrowane podczas wczytywania.
    \item Generator umożliwia tworzenie pytań wielokrotnego wyboru (cztery odpowiedzi, w tym minimum jedna poprawna).
    \item Aplikacja rozwiązująca pozwala na wczytanie pliku, odliczanie czasu za pomocą zegara oraz automatyczne obliczenie punktów po zakończeniu.
\end{itemize}

\subsection{Stos technologiczny}
\begin{itemize}
    \item \textbf{Język programowania:} C\#
    \item \textbf{Platforma:} .NET
    \item \textbf{Interfejs graficzny:} Windows Forms
\end{itemize}

\subsection{Oczekiwane rezultaty projektu}
Gotowy system pozwalający na bezproblemowe tworzenie szyfrowanych zestawów pytań, a następnie ich bezpieczne wczytywanie w aplikacji rozwiązującej z przyjaznym interfejsem zapobiegającym błędom użytkownika.

\section{Realizacja projektu}

\subsection{Moduł Rozwiązywania Quizu (\authorA)}

Moduł \texttt{QuizSolver\_M} został zrealizowany zgodnie z wzorcem \textbf{MVP (Model-View-Presenter)},
który zapewnia separację logiki biznesowej od interfejsu użytkownika.
Projekt korzysta ze współdzielonych bibliotek \texttt{QuizModel} i \texttt{QuizServices},
zawierających klasy modelu (\texttt{Quiz}, \texttt{Question}, \texttt{Answer})
oraz serwisy serializacji i szyfrowania.

\subsubsection{Architektura MVP}

\begin{itemize}
    \item \textbf{Model} --- klasy \texttt{Quiz}, \texttt{Question}, \texttt{Answer}
          oraz \texttt{QuizSerializer} i \texttt{EncryptionService} współdzielone z generatorem.
    \item \textbf{View (\texttt{SolverView})} --- formularz WinForms implementujący interfejs
          \texttt{ISolverView}. Widok nie zawiera logiki --- wystawia zdarzenia i aktualizuje UI
          na polecenie prezentera.
    \item \textbf{Presenter (\texttt{SolverPresenter})} --- subskrybuje zdarzenia widoku,
          steruje przebiegiem quizu i aktualizuje widok przez metody interfejsu.
\end{itemize}

\subsubsection{Wczytywanie i automatyczne deszyfrowanie}

Po wybraniu pliku przez użytkownika (\texttt{OpenFileDialog}) prezenter próbuje
odszyfrować go kaskadowo, testując kolejno strategie \textbf{AES}, \textbf{TripleDES}
i \textbf{Plaintext}. Dzięki temu aplikacja jest kompatybilna z każdym trybem zapisu
generatora bez konieczności ręcznego wyboru algorytmu.

\begin{lstlisting}
var strategies = new List<IEncryptionStrategy>
{
    EncryptionService.Instance.GetAesStrategy(),
    EncryptionService.Instance.GetTripleDesStrategy(),
    new PlaintextStrategy()
};
foreach (var strategy in strategies)
{
    try
    {
        var serializer = new QuizSerializer(strategy);
        _quiz = serializer.Deserialize(dialog.FileName);
        _questions = _quiz.Questions.Items.Values.ToList();
        _userAnswers = _questions.ToDictionary(
            q => q.Id,
            q => new List<bool> { false, false, false, false });
        _view.SetQuizTitle(_quiz.Name, _questions.Count);
        _view.EnableStart(true);
        return;
    }
    catch { }
}
_view.DisplayError("Nie udalo sie wczytac quizu.");
\end{lstlisting}

\subsubsection{Sterowanie przebiegiem quizu i timer}

Prezenter kontroluje stan interfejsu poprzez włączanie i wyłączanie elementów UI:
\begin{itemize}
    \item \textbf{Po wczytaniu pliku:} aktywowany jest przycisk „Start".
    \item \textbf{Po kliknięciu „Start":} uruchamiany jest \texttt{Timer} (interwał 1000\,ms),
          aktywowane są przyciski nawigacji i „Zakończ".
    \item \textbf{Po kliknięciu „Zakończ":} timer zatrzymywany, checkboxy blokowane,
          obliczany wynik i wyświetlane kolorowe oznaczenia odpowiedzi.
\end{itemize}

\begin{lstlisting}
_timer.Interval = 1000;
_timer.Tick += (s, e) =>
    _view.UpdateTimer((DateTime.Now - _startTime).ToString(@"mm\:ss"));
\end{lstlisting}

\subsubsection{Punktacja i weryfikacja odpowiedzi}

Odpowiedzi użytkownika przechowywane są w słowniku \texttt{\_userAnswers} (klucz: \texttt{Guid} pytania,
wartość: lista 4 wartości \texttt{bool}). Użytkownik otrzymuje 1 punkt za pytanie tylko wtedy,
gdy jego zaznaczenia są identyczne z zestawem poprawnych odpowiedzi (metoda \texttt{SequenceEqual}).

\begin{lstlisting}
private int CalculateScore()
{
    int score = 0;
    foreach (var q in _questions)
    {
        var correct = q.Answers.Items.Select(a => a.IsCorrect).ToList();
        var user = _userAnswers[q.Id];
        if (correct.SequenceEqual(user)) score++;
    }
    return score;
}
\end{lstlisting}

\subsubsection{Podgląd poprawnych i błędnych odpowiedzi}

Po zakończeniu quizu użytkownik może przeglądać pytania z oznaczeniem kolorystycznym odpowiedzi:
\begin{itemize}
    \item \textcolor{green!60!black}{\textbf{Zielony}} --- zaznaczona i poprawna,
    \item \textcolor{orange}{\textbf{Pomarańczowy}} --- poprawna, ale niezaznaczona (pominięta),
    \item \textcolor{red}{\textbf{Czerwony}} --- zaznaczona, ale błędna.
\end{itemize}

Checkboxy pozostają wizualnie włączone (WinForms przy \texttt{Enabled = false} nadpisuje
\texttt{ForeColor} na szary, co uniemożliwiałoby wyświetlenie kolorów), a blokada zmian
realizowana jest przez flagę \texttt{\_quizEnded} w presenterze.

\begin{lstlisting}
private void OnEndQuiz()
{
    _timer.Stop();
    _elapsed = DateTime.Now - _startTime;
    _quizEnded = true;
    _view.EnableEnd(false);
    _view.EnableNavigation(true);
    int score = CalculateScore();
    ShowCorrectAnswers();
    _view.ShowResult(score, _questions.Count, _elapsed);
}
\end{lstlisting}


% ====
% MIEJSCE NA GENERATOR QUIZU - OSOBA 2
% ====
\subsection{Moduł Generatora Quizu (\authorB)}
\textit{\todo{Tutaj należy opisać realizację generatora quizu. Jak zrealizowano tworzenie pytań, mechanizm szyfrowania i zapisu do pliku dyskowego.}}

% [Wklej kod lub opis tutaj]


% ====
% MIEJSCE NA DRUGIE ROZWIĄZANIE QUIZU - OSOBA 3
% ====
\subsection{Alternatywne Rozwiązanie Quizu / Moduł 3 (\authorC)}
\textit{\todo{Tutaj należy opisać drugą aplikację do rozwiązywania quizu lub rozszerzenie jej o dodatkowe funkcje (np. inna architektura WPF/MVVM).}}

% [Wklej kod lub opis tutaj]


\section{Wnioski}
\begin{itemize}
    \item \textbf{Spostrzeżenia:} Zastosowanie wzorca projektowego MVP w aplikacji Windows Forms znacząco ułatwiło separację logiki od interfejsu graficznego.
    \item \textbf{Osiągnięcia:} Zrealizowano poprawnie funkcjonującą warstwę prezentera oraz szyfrowanie/deszyfrowanie plików projektowych.
    \item \textbf{Potencjał rozwoju:} Aplikację można w przyszłości rozbudować o możliwość eksportu statystyk do bazy danych, czy obsługę multimediów w pytaniach.
\end{itemize}

\section{Bibliografia}
\begin{enumerate}
    \item Wytyczne do zadania nr 3 - ,,Quiz'', semestr 4, kierunek Informatyka.
    \item Dokumentacja języka C\# i platformy .NET. Microsoft Docs.
\end{enumerate}

\end{document}